\section{输入、公式与输出：每一个前瞻数字都有口径}

本附录列出主模型的核心输入。除公司预告区间和历史财务外，2026--2027E 均为 House 假设。模型不用“设计产能×行业价格”替代实际经营；锂业务收入为销量乘以实现收入代理，锂毛利为收入乘毛利率，归母由合并毛利扣除期间费用、税和少数股东损益后的情景输出。

\begin{exhibitbox}[表 A-4：2026E House 核心假设]
\begin{tabularx}{\exhibitboxwidth}{L{3.6cm}R{2.0cm}R{2.0cm}R{2.0cm}X}
\toprule
输入 & 熊 & 基准 & 牛 & 口径 \\
\midrule
H1 归母净利（亿元） & 11.0 & 12.0 & 13.0 & 公司预告下/中/上端 \\
H2 归母净利（亿元） & 4.0 & 8.1 & 11.7 & House 独立预测 \\
锂盐销量（吨） & 75,000 & 90,000 & 100,000 & 不按设计产能 \\
实现收入代理（元/吨） & 95,000 & 105,000 & 112,000 & 不等于现货价或实际 ASP \\
锂业务毛利率 & 25\% & 32\% & 35\% & 不给未证实成本溢价 \\
民爆产品收入 / 毛利率 & 21.5 / 42\% & 22.0 / 44\% & 23.0 / 46\% & 亿元 / 百分比 \\
爆破服务收入 / 毛利率 & 12.0 / 18\% & 12.5 / 19\% & 13.0 / 20\% & 亿元 / 百分比 \\
OCF / 资本开支 / FCF（亿元） & -1.0 / 3.0 / -4.0 & 2.0 / 3.5 / -1.5 & 6.0 / 4.0 / 2.0 & 现金另行验证 \\
\bottomrule
\end{tabularx}
\end{exhibitbox}

\section{估值参数与敏感性：不让单一倍数主导结论}

每个情景的主估值由 70\% 的周期归一化 PE 和 30\% 的 PB/ROE 交叉检查构成。PE 法又按 70\% 的 2026E EPS 与 30\% 的 2027E EPS 加权，避免将 2026 利润高点永久化。概率为熊/基/牛 30/50/20；最终多锚目标将概率加权内在价值赋 85\% 权重、市场锚 15\%，外部目标 0\%。

\begin{exhibitbox}[表 A-5：估值参数、公式与敏感性]
\begin{tabularx}{\exhibitboxwidth}{L{3.5cm}R{2.0cm}R{2.0cm}R{2.0cm}X}
\toprule
项目 & 熊 & 基准 & 牛 & 说明 \\
\midrule
2026E EPS & 1.301 & 1.744 & 2.143 & 归母 ÷ 11.5256 亿股 \\
2027E EPS & 1.015 & 1.237 & 1.614 & 周期归一化输入 \\
2026E / 2027E PE & 8.0x / 9.0x & 10.5x / 12.0x & 13.0x / 14.0x & 周期校准 \\
PE 法价值 & 10.03 元 & 17.27 元 & 26.28 元 & 70\%×2026E + 30\%×2027E \\
PB/ROE 价值 & 10.79 元 & 17.42 元 & 25.60 元 & BPS 与 PB 交叉检查 \\
情景内在价值 & 10.26 元 & 17.32 元 & 26.07 元 & 70\% PE + 30\% PB/ROE \\
情景概率 & 30\% & 50\% & 20\% & 现金与 H2 证据折价 \\
\bottomrule
\end{tabularx}
\end{exhibitbox}

\section{可复算性与模型限制}

概率加权内在价值为 $10.26\times30\%+17.32\times50\%+26.07\times20\%=16.95$ 元；最终值为 $16.95\times85\%+17.75\times15\%=17.07$ 元，展示为 17.0 元。模型不包含实际自供比例、实际单位成本、套保损益、低价库存收益、客户订单、民爆新项目和硫化锂商业化；这些项目的正确处理是零信用，而不是以模糊的“潜在弹性”写入 EPS。

\section{单变量敏感性：先识别变量，再谈上行}

下表不是对单项变化的精确利润弹性承诺，而是展示基准情景最敏感的变量及其传导方向。收入代理、销量和毛利率必须同时受正式 H1 的收入、毛利和营运资本约束；因此不能把一项行业价格变动孤立地转化为目标价。

\begin{exhibitbox}[表 A-6：基准情景的变量敏感性与验证要求]
\begin{tabularx}{\exhibitboxwidth}{L{3.0cm}L{2.7cm}X X}
\toprule
变量 & 基准输入 & 主要传导 & 能够改变假设的证据 \\
\midrule
锂盐销量 & 9.0 万吨 & 直接改变锂业务收入与规模效应 & H1 销量或等效运营数据 \\
实现收入代理 & 10.5 万元/吨 & 改变收入，但不等同现货价格 & 分部收入、产品结构及结算口径 \\
锂业务毛利率 & 32\% & 同时反映成本、库存和效率 & 分部毛利、成本和存货说明 \\
H2 归母 & 8.1 亿元 & 直接决定全年 EPS 与估值中枢 & H1 桥与下半年运营证据 \\
经营现金流 & 2.0 亿元 & 影响 PE 折价和牛市概率 & 现金流量表、应收和存货变动 \\
民爆产品毛利率 & 44\% & 影响防御性利润池质量 & 分部成本、项目执行和回款 \\
2027E EPS & 1.237 元 & 限制高点 PE 外推 & 2026H2 单位经济及周期条件 \\
\bottomrule
\end{tabularx}
\end{exhibitbox}

一个实用的研究纪律是：只有当新的公司披露同时改善至少一项经营变量和至少一项现金/资产负债表变量时，才提高牛情景概率。只改善商品价格或只改善预告标题，均不足以改变 17.0 元目标的估值权重。

\sourcenote{机器可读估值和增长模型；复算审计见模型审计工作底稿。}
